\documentclass[11pt]{article}

\usepackage[a4paper,margin=1in]{geometry}
\usepackage{amsmath,amssymb}
\usepackage{bm}
\usepackage{booktabs}
\usepackage{siunitx}
\usepackage{hyperref}
\usepackage{enumitem}

\hypersetup{colorlinks=true,linkcolor=blue,urlcolor=blue}

\newcommand{\R}{\mathbb{R}}
\newcommand{\SOO}{\mathrm{SO}(3)}
\newcommand{\SEE}{\mathrm{SE}(3)}
\newcommand{\Rx}{R_x}
\newcommand{\Ry}{R_y}
\newcommand{\Rz}{R_z}
\newcommand{\skewm}[1]{\left[#1\right]_{\times}}

\title{Forward and Inverse Kinematics of the \texttt{new\_arm} Manipulator}
\author{Robot description: \texttt{new\_arm\_description} \\
        Solver package: \texttt{new\_arm\_ik}}
\date{2026-05-31}

\begin{document}
\maketitle

\begin{abstract}
This document describes the kinematics of the \texttt{new\_arm}, a 4-DOF
revolute manipulator, as implemented in the \texttt{new\_arm\_ik} ROS\,2
package. The forward kinematics is built directly from the robot's URDF so
that the solver and the model rendered in RViz are identical by construction.
The inverse kinematics is formulated as a bounded nonlinear least-squares
problem with a continuity-preserving regularizer and a random-restart
fallback. We also document a tuning fix that reduced the median position
residual from \SI{20}{\milli\meter} to below \SI{1}{\micro\meter} on
reachable targets.
\end{abstract}

\tableofcontents

\section{Robot model}

The \texttt{new\_arm} is a serial chain of four revolute joints connecting
five links:
\[
\texttt{arm\_base} \;\xrightarrow{\;q_1\;}\;
\texttt{after\_base\_full} \;\xrightarrow{\;q_2\;}\;
\texttt{link1} \;\xrightarrow{\;q_3\;}\;
\texttt{link2} \;\xrightarrow{\;q_4\;}\;
\texttt{end\_effector}.
\]
The kinematic parameters are read verbatim from
\texttt{new\_arm\_description/urdf/new\_arm.urdf}. Each joint $i$ contributes
a fixed origin transform (a translation $\bm{p}_i\in\R^3$ and a fixed
orientation given by roll--pitch--yaw angles $\bm{r}_i=(\phi_i,\theta_i,\psi_i)$)
followed by a rotation of $q_i$ about the joint axis $\hat{\bm{a}}_i\in\R^3$.
Table~\ref{tab:joints} lists the values.

\begin{table}[h]
\centering
\caption{Joint parameters from the URDF. Origins $\bm{p}_i$ are in metres;
all fixed RPY orientations are zero. Limits are symmetric.}
\label{tab:joints}
\small
\begin{tabular}{llrrr}
\toprule
Joint & Axis $\hat{\bm{a}}_i$ & Origin $\bm{p}_i=(x,y,z)$ \si{\meter} & Limit \si{\radian} & Effort \\
\midrule
\texttt{after\_base\_full\_joint} & $(0,0,1)$ & $(0.00043,\,-0.00184,\,0.08396)$ & $[-3.14,\,3.14]$ & 10.0 \\
\texttt{link1\_joint}             & $(0,1,0)$ & $(0.00814,\,-0.00043,\,0.00162)$ & $[-3.14,\,3.14]$ & 10.0 \\
\texttt{link2\_joint}             & $(0,1,0)$ & $(0.00100,\,-0.00856,\,0.12840)$ & $[-3.14,\,3.14]$ & 10.0 \\
\texttt{end\_effector\_joint}     & $(0,1,0)$ & $(0.12532,\,0.01099,\,-0.05778)$ & $[-3.14,\,3.14]$ & 5.0  \\
\bottomrule
\end{tabular}
\end{table}

The first joint rotates about $z$ (a base yaw); the remaining three rotate
about $y$ (pitches). The reachable workspace is a region of radius
$\approx \SI{0.25}{\meter}$. In the current interface the fourth joint
$q_4$ (\texttt{end\_effector\_joint}) is held at $0$; the position IK has
three Cartesian targets and four joint variables, so the arm is kinematically
redundant by one degree of freedom.

\section{Rotation primitives}

\subsection{Roll--pitch--yaw to a rotation matrix}

URDF \texttt{origin} orientations use the ROS fixed-axis roll--pitch--yaw
convention. With roll $\phi$ about $x$, pitch $\theta$ about $y$, and yaw
$\psi$ about $z$, the composed rotation is the intrinsic $z\,y\,x$ product
\begin{equation}
R(\phi,\theta,\psi) \;=\; \Rz(\psi)\,\Ry(\theta)\,\Rx(\phi),
\end{equation}
where
\[
\Rx(\phi)=\begin{bmatrix}1&0&0\\0&c_\phi&-s_\phi\\0&s_\phi&c_\phi\end{bmatrix},\quad
\Ry(\theta)=\begin{bmatrix}c_\theta&0&s_\theta\\0&1&0\\-s_\theta&0&c_\theta\end{bmatrix},\quad
\Rz(\psi)=\begin{bmatrix}c_\psi&-s_\psi&0\\s_\psi&c_\psi&0\\0&0&1\end{bmatrix},
\]
with the shorthand $c_\alpha=\cos\alpha$, $s_\alpha=\sin\alpha$. Expanding the
product gives the closed form implemented in \texttt{rpy\_to\_matrix}:
\begin{equation}
R(\phi,\theta,\psi)=
\begin{bmatrix}
c_\psi c_\theta & c_\psi s_\theta s_\phi - s_\psi c_\phi & c_\psi s_\theta c_\phi + s_\psi s_\phi \\
s_\psi c_\theta & s_\psi s_\theta s_\phi + c_\psi c_\phi & s_\psi s_\theta c_\phi - c_\psi s_\phi \\
-s_\theta       & c_\theta s_\phi                        & c_\theta c_\phi
\end{bmatrix}.
\end{equation}

\subsection{Axis--angle (Rodrigues) rotation}

Each actuated joint rotates by $q_i$ about a unit axis
$\hat{\bm{a}}=(a_x,a_y,a_z)$. The rotation matrix is given by Rodrigues'
formula
\begin{equation}
R(\hat{\bm{a}},q)=I + (\sin q)\,\skewm{\hat{\bm{a}}}
        + (1-\cos q)\,\skewm{\hat{\bm{a}}}^2,
\end{equation}
where $\skewm{\hat{\bm{a}}}$ is the skew-symmetric cross-product matrix
\[
\skewm{\hat{\bm{a}}}=
\begin{bmatrix}0&-a_z&a_y\\ a_z&0&-a_x\\ -a_y&a_x&0\end{bmatrix}.
\]
Written componentwise (the form in \texttt{axis\_angle\_matrix}), with
$c=\cos q$, $s=\sin q$, and $C=1-c$:
\begin{equation}
R(\hat{\bm{a}},q)=
\begin{bmatrix}
c+a_x^2 C        & a_x a_y C - a_z s & a_x a_z C + a_y s \\
a_y a_x C + a_z s & c+a_y^2 C        & a_y a_z C - a_x s \\
a_z a_x C - a_y s & a_z a_y C + a_x s & c+a_z^2 C
\end{bmatrix}.
\end{equation}

\subsection{Homogeneous transforms}

A rotation $R\in\SOO$ and translation $\bm{t}\in\R^3$ are packed into a
homogeneous transform $T\in\SEE$:
\begin{equation}
T(R,\bm{t})=\begin{bmatrix}R & \bm{t}\\ \bm{0}^\top & 1\end{bmatrix}\in\R^{4\times4}.
\end{equation}

\section{Forward kinematics}

The pose of the end-effector relative to the base is the ordered product of,
for each joint, the fixed origin transform followed by the variable joint
rotation:
\begin{equation}
\boxed{\;
T^{\text{base}}_{\text{ee}}(\bm{q})
  = \prod_{i=1}^{4}
      \underbrace{T\!\big(R(\bm{r}_i),\,\bm{p}_i\big)}_{\text{fixed origin}}
      \;\underbrace{T\!\big(R(\hat{\bm{a}}_i, q_i),\,\bm{0}\big)}_{\text{joint }q_i}
\;}
\end{equation}
where $\bm{q}=(q_1,q_2,q_3,q_4)$. This is exactly the loop in
\texttt{Chain.fk}. The end-effector position used by the IK is the
translation block
\begin{equation}
\bm{x}(\bm{q}) \;=\; \big[T^{\text{base}}_{\text{ee}}(\bm{q})\big]_{1:3,\,4}\;\in\;\R^3 .
\end{equation}

Because the parameters $\{\bm{p}_i,\bm{r}_i,\hat{\bm{a}}_i\}$ are parsed
straight from the URDF, $\bm{x}(\bm{q})$ coincides with the frame that
\texttt{robot\_state\_publisher} broadcasts to RViz; there is no separate
hand-derived model to drift out of sync.

As a reference, the home pose $\bm{q}=\bm{0}$ gives
\[
\bm{x}(\bm{0}) \approx (0.1349,\;0.0002,\;0.1562)\ \si{\meter}.
\]

\section{Inverse kinematics}

\subsection{Problem statement}

Given a Cartesian target $\bm{x}^\star\in\R^3$, find joint angles
$\bm{q}\in\R^4$ within the box limits $[\bm{q}_{\min},\bm{q}_{\max}]$ whose
forward kinematics reaches the target. Since the map $\bm{x}(\bm{q})$ is
nonlinear and the arm is redundant, we solve a regularized nonlinear
least-squares problem rather than a closed-form inverse.

\subsection{Residual and cost}

Let $\bm{q}_{\text{ref}}$ be a reference configuration (the current pose, for
continuity). Define the stacked residual $\bm{r}:\R^4\to\R^{7}$
\begin{equation}
\bm{r}(\bm{q})=
\begin{bmatrix}
w_p\,\big(\bm{x}(\bm{q})-\bm{x}^\star\big)\\[2pt]
\lambda\,\big(\bm{q}-\bm{q}_{\text{ref}}\big)
\end{bmatrix}
\in\R^{3+4},
\end{equation}
combining a \emph{position} term (weight $w_p$, three rows) and a
\emph{regularization} term (weight $\lambda$, four rows). The solver
minimizes
\begin{equation}
\bm{q}^\star=\arg\min_{\bm{q}_{\min}\le\bm{q}\le\bm{q}_{\max}}
   \tfrac{1}{2}\,\lVert \bm{r}(\bm{q})\rVert_2^2
 = \arg\min_{\bm{q}}\ \tfrac{1}{2}\Big(
      w_p^2\,\lVert \bm{x}(\bm{q})-\bm{x}^\star\rVert^2
    + \lambda^2\,\lVert \bm{q}-\bm{q}_{\text{ref}}\rVert^2 \Big).
\end{equation}
The position term drives the end-effector to the target. The regularization
term has two jobs: it resolves the one-dimensional redundancy by selecting,
among all exact solutions, the one nearest $\bm{q}_{\text{ref}}$, and it keeps
successive solves close together so the arm moves smoothly in RViz with no
joint flips.

This is solved with \texttt{scipy.optimize.least\_squares} (Trust Region
Reflective), which respects the box bounds and uses a finite-difference
Jacobian $J=\partial\bm{r}/\partial\bm{q}$.

\subsection{The bias--accuracy trade-off in \texorpdfstring{$\lambda$}{lambda}}

At a solution the two objectives compete: increasing $\lambda$ pulls
$\bm{q}^\star$ toward $\bm{q}_{\text{ref}}$ at the expense of position
accuracy. Concretely, near a minimum the position error scales like
\begin{equation}
\lVert \bm{x}(\bm{q}^\star)-\bm{x}^\star\rVert
   \;\sim\; \frac{\lambda^2}{w_p^2}\,\lVert J_x^{+}\rVert\,
            \lVert \bm{q}^\star-\bm{q}_{\text{ref}}\rVert,
\end{equation}
where $J_x$ is the $3\times4$ position Jacobian and $J_x^{+}$ its pseudoinverse.
The residual therefore grows roughly quadratically in $\lambda$. This is the
mechanism behind the bug described in Section~\ref{sec:fix}: an
over-large $\lambda$ produced centimetre-scale errors on every target.

\subsection{Random-restart fallback}

Even with $\lambda\to 0$, a single local optimization seeded from one point
can stall in a local minimum (the position objective is non-convex), which
empirically happened on about \SI{2}{\percent} of reachable targets, with
residuals up to $\sim\SI{100}{\milli\meter}$. The solver therefore uses a
two-phase strategy:

\begin{enumerate}[leftmargin=2em]
  \item \textbf{Continuity solve.} Seed at $\bm{q}_{\text{init}}=\bm{q}_{\text{ref}}=$
        the current pose and solve once. If the resulting residual is within
        a tolerance $\varepsilon$ (default \SI{0.5}{\milli\meter}), accept it.
        This is the fast common path and yields smooth motion.
  \item \textbf{Restart fallback.} Otherwise, draw up to $N$ random seeds
        uniformly from the joint box,
        $\bm{q}^{(k)}\sim\mathcal{U}(\bm{q}_{\min},\bm{q}_{\max})$, solve from
        each (regularizing toward that seed, i.e.\ no continuity bias), and
        keep the lowest-residual result. Stop early as soon as the tolerance
        is met.
\end{enumerate}

Formally, the returned solution is
\begin{equation}
\bm{q}^\star=
\arg\min_{\bm{q}\in\{\,\bm{q}^{(0)},\dots,\bm{q}^{(N)}\,\}}
   \big\lVert \bm{x}(\bm{q})-\bm{x}^\star\big\rVert,
\end{equation}
where $\bm{q}^{(0)}$ is the continuity solve and
$\bm{q}^{(1)},\dots,\bm{q}^{(N)}$ are the restart solves (evaluated only when
needed). The hard targets pay the extra cost; the easy ones do not.

\section{Solver parameters}

\begin{table}[h]
\centering
\caption{\texttt{IKSolver} parameters and their meaning.}
\label{tab:params}
\small
\begin{tabular}{lll}
\toprule
Symbol & Code name & Default \\
\midrule
$w_p$         & \texttt{position\_weight} & $1.0$ \\
$\lambda$     & \texttt{regularization}   & $10^{-4}$ \\
$\varepsilon$ & \texttt{tolerance}        & $5\times10^{-4}\ \si{\meter}$ \\
$N$           & \texttt{restarts}         & $6$ \\
--            & \texttt{max\_iter}        & $200$ \\
\bottomrule
\end{tabular}
\end{table}

\section{The tuning fix}
\label{sec:fix}

The original solver shipped with $\lambda=0.05$ and no restart fallback. By
the relation in the trade-off subsection, that regularization weight was large
enough to dominate the position objective, biasing every solution toward the
seed. The effect was measured by generating guaranteed-reachable targets
(sample random $\bm{q}_{\text{true}}$, compute $\bm{x}(\bm{q}_{\text{true}})$
by FK, then ask the IK to recover the target); this isolates solver quality
from workspace limits.

\begin{table}[h]
\centering
\caption{Position residual on reachable targets, before and after. Lower is
better; ``miss rate'' is the fraction of targets with residual
$>\SI{1}{\milli\meter}$.}
\label{tab:results}
\small
\begin{tabular}{lrrr}
\toprule
Configuration & Median & Worst case & Miss rate \\
\midrule
$\lambda=0.05$ (original)              & \SI{20}{\milli\meter}   & \SI{46}{\milli\meter} & \SI{100}{\percent} \\
$\lambda=0.005$                       & \SI{0.36}{\milli\meter} & \SI{3.1}{\milli\meter}& \SI{12}{\percent} \\
$\lambda=0$, single start             & \SI{0}{\milli\meter}    & \SI{108}{\milli\meter}& \SI{2}{\percent} \\
$\lambda=10^{-4}$ + restart (current) & $<\SI{1}{\micro\meter}$ & \SI{0.022}{\milli\meter}& \SI{0}{\percent} \\
\bottomrule
\end{tabular}
\end{table}

The adopted configuration ($\lambda=10^{-4}$ with the two-phase restart
fallback) achieves sub-micrometre median accuracy and zero misses over $500$
reachable targets, while preserving smooth motion: stepping the target along a
line produced a maximum single-step joint change of only $0.68^\circ$.

\section{ROS\,2 integration}

The solver is wrapped by \texttt{ik\_gui\_node.py}, a node that:
\begin{itemize}[leftmargin=2em]
  \item parses the chain once with \texttt{parse\_urdf\_chain(urdf, root, tip)};
  \item exposes a Tkinter panel for entering an $(x,y,z)$ target in metres;
  \item calls \texttt{solver.solve(target, q\_init=self.current\_q)} on demand,
        feeding the current pose as the continuity reference;
  \item publishes the resulting joint vector on \texttt{/joint\_states} at
        \SI{20}{\hertz}, which \texttt{robot\_state\_publisher} turns into TF
        frames for RViz.
\end{itemize}
No change to the node was required for the tuning fix; it benefits
automatically because the solver's public \texttt{solve} signature is
unchanged.

\paragraph{Running it.}
\begin{verbatim}
colcon build --packages-select new_arm_description new_arm_ik
source install/setup.bash
ros2 launch new_arm_ik ik_interface.launch.py
\end{verbatim}

\end{document}
